% NAEGAUGEHIGHGIRTH1 2026-10-17 1800 Asia/Seoul
\documentclass[11pt]{article}
\usepackage[T1]{fontenc}
\usepackage{lmodern,amsmath,amssymb,amsthm,booktabs,hyperref}
\usepackage[margin=0.9in]{geometry}
\newtheorem{theorem}{Theorem}
\newtheorem{lemma}[theorem]{Lemma}
\newtheorem{corollary}[theorem]{Corollary}
\newcommand{\EXACT}{\textbf{PROVED/EXACT}}
\newcommand{\KNOWN}{\textbf{KNOWN THEOREM}}
\newcommand{\COMP}{\textbf{COMPUTATIONAL}}
\newcommand{\CAND}{\textbf{CANDIDATE}}
\newcommand{\OPEN}{\textbf{OPEN}}
\newcommand{\RET}{\textbf{FALSIFIED/RETIRED}}
\title{NAEGAUGEHIGHGIRTH1:\\Gauge-Code Independence and Logarithmic Moment Matching for Exact-Density NAE-UNSAT}
\author{P versus NP project}
\date{17 October 2026}
\begin{document}\maketitle

\begin{center}\fbox{\parbox{0.94\linewidth}{
\textbf{Status.} The gauge-rank/cycle-defect theorem, forest independence theorem, first-moment UNSAT signing lemma, and diluted-pattern moment bound are \EXACT at their stated models, pending independent review. Ellis--Linial high-girth constructions, Kun sparse incomparability, Darmann--D\"ocker bounded-occurrence NAE hardness, and the 2026 derandomized robust-sunflower theorem are \KNOWN. The logarithmic moment-matching distribution is \EXACT conditional only on the cited Ellis--Linial explicit support sequence. The circuit-to-bounded-moment approximation theorem remains \OPEN, so $P\ne NP$ remains \OPEN.
}}\end{center}

\section{IDEAS PROPOSED BEFORE LITERATURE}
\paragraph{A: fixed-girth hardness plus gauge randomization.} Use Kun's sparse-incomparability theorem to separate the two requirements that previously conflicted: computational hardness and locally tree-like support. Then random variable complementation should erase signed local structure on every forest.

\paragraph{B: cycle-defect code.} Treat the four signed NAE classes on each support as two bits. The map from vertex gauge bits to edge-sign bits should have an exact rank defect equal to the incidence-cycle rank. Missing theorem: prove this for every 3-uniform support, not just trees.

\paragraph{C: logarithmic high-girth UNSAT cores.} Instead of requiring the hard family itself to have growing girth, use a bounded-degree high-girth support and choose one UNSAT signing by the first moment. Dilute a sublinear core into the exact-density-$4/3$ NAE slice. Missing theorem: uniform low-order clause-moment matching after embedding, gauge randomization, and filler.

\paragraph{D: exponent-indexed locality.} For contradiction against a hypothetical polynomial circuit exponent $K$, it is enough for the positive distribution to be pseudorandom through the locality order required for that fixed $K$. A single universally optimal parameter tuple is unnecessary.

\section{LITERATURE CHECK AND REVISIONS}
Kun proves that every CSP is polynomial-time equivalent to its restriction to structures of arbitrarily large fixed girth. His Theorem 2 constructs, for every fixed girth parameter, a structure of larger girth preserving homomorphism to every fixed-size target. The paper explicitly states that the theorem applies to uniform hypergraphs. Applied to the two-element NAE template, this makes NAE satisfiability hard even after any fixed girth requirement. The proof also uses bounded-degree expander structures, so starting from the degree-four hard NAE source keeps degree bounded for each fixed target girth.

Darmann--D\"ocker prove Positive NAE-3-SAT NP-complete for linear formulas in which every variable occurs exactly four times. This remains the clean hard source, but this run does not need the lifted family to preserve exact degree four.

Ellis--Linial give explicit logarithmic-girth regular hypergraphs when the degree is a multiple of the uniformity. For $r=3,d=9$, their Ramanujan-star construction gives an infinite sequence of linear 3-uniform 9-regular supports on $v$ vertices with
\[g(H)\ge \gamma\log v-O(1),\qquad \gamma=\frac{2}{3\log 11}.\]
The choice $d=9$ is also above the exact first-moment threshold for a random signed NAE formula to be UNSAT with positive probability.

Cavalar--Fabris--Mukhopadhyay--Srinivasan--Yehudayoff show in 2026 that robust-sunflower statements can be derandomized to $O(w)$-wise independent distributions, with further locally-independent approximately-identical marginal variants. Their abstract also notes the older $O(w^2)$-wise route through DNF fooling. This makes exact local independence of NAE signs directly relevant, but it does not itself approximate an arbitrary monotone circuit on our same-layer distributions.

\section{CONSTRUCTION OR PROOF ATTEMPT}
\subsection{Gauge code and exact cycle defect}
Let $H=(V,E)$ be a 3-uniform support hypergraph. A signed NAE predicate on $e=\{u,v,w\}$ is represented by a vector in
\[Q_e=\mathbb F_2^3/\langle(1,1,1)\rangle\cong\mathbb F_2^2,\]
because globally complementing all three literals does not change NAE. Fix any base signing $\sigma=(\sigma_e)_{e\in E}$. For a vertex gauge $x\in\mathbb F_2^V$, define
\[\lambda_e(x)=\sigma_e+x|_e\pmod{\langle111\rangle}.\]
Complementing variable $v$ exactly adds $x_v$ to every incident literal, so every gauge transform is satisfiability-equivalent to the base formula.

\begin{theorem}[Gauge-rank/cycle-defect theorem]
Let $J\subseteq E$ have $m$ edges, use $v_J$ vertices, and have $c_J$ connected components after isolated vertices are discarded. The linear gauge map
\[\delta_J:\mathbb F_2^{V(J)}\to\bigoplus_{e\in J}Q_e\]
has
\[\boxed{\operatorname{rank}\delta_J=v_J-c_J.}\]
Hence its codimension in the $2m$-dimensional edge-label space is
\[\boxed{\xi(J)=2m-v_J+c_J.}\]
This equals the cyclomatic number of the bipartite incidence graph of $J$.
\end{theorem}
\begin{proof}
A vector lies in $\ker\delta_J$ exactly when its three values are equal on every hyperedge. By connectedness this means the vector is constant on each connected component, so the kernel dimension is $c_J$. Rank-nullity gives $v_J-c_J$. The incidence graph has $v_J+m$ vertices, $3m$ incidence edges and $c_J$ components, so its cycle rank is $3m-(v_J+m)+c_J=2m-v_J+c_J$.
\end{proof}

\begin{corollary}[Forest independence]
If $J$ is a Berge hyperforest then $v_J=2m+c_J$, so $\xi(J)=0$ and $\delta_J$ is onto. Therefore, under a uniform random vertex gauge, the four NAE sign classes on the edges of $J$ are exactly mutually independent and uniform.
More generally, for arbitrary $J$, exactly a $2^{-\xi(J)}$ fraction of all sign-label vectors are attainable, and each attainable vector has probability
\[4^{-m}2^{\xi(J)}.\]
\end{corollary}

Thus a support of girth $g$ has exactly $(g-1)$-wise independent signed edge labels under the gauge distribution. Equivalently the associated linear edge-label code has dual edge-support distance at least $g$.

\subsection{High-girth UNSAT signed NAE cores}
Let $H$ be a 3-uniform 9-regular support on $v$ vertices. It has $k=3v$ edges. Independently choose one of the four signed NAE predicates on each edge. For any fixed assignment, exactly one of the four sign classes is falsified, so
\[\mathbb E[\#\text{satisfying assignments}]=2^v(3/4)^{3v}=(27/32)^v<1.\]
Since the number of satisfying assignments is a nonnegative integer, some signing is UNSAT.

\begin{corollary}[Logarithmic-girth UNSAT sequence]
Along the Ellis--Linial 9-regular sequence there are signed NAE-3CNFs $K_v$ with
\[|E(K_v)|=3v,\qquad K_v\text{ UNSAT},\qquad g(K_v)\ge\frac{2}{3\log11}\log v-O(1).\]
Random gauge transforms of $K_v$ remain UNSAT and have exactly independent uniform sign labels on every edge set below the girth.
\end{corollary}

\subsection{Dilution into the exact-density NAE slice}
Let the ambient signed NAE coordinate universe on $q$ variables have
\[M=4\binom q3,\qquad t=4q/3\quad(3\mid q).\]
Choose $v=q^{\alpha+o(1)}$ along the core subsequence, with $0<\alpha<1$. Randomly inject the core vertices into $[q]$, apply a uniform vertex gauge, and then choose $t-k$ filler clause coordinates uniformly from the $M-k$ coordinates not already occupied by the core. Call the resulting exact-$t$ UNSAT formula $Y$.

For a fixed set $A$ of $j$ distinct signed clause coordinates, let
\[u_j=\frac{(t)_j}{(M)_j}\]
be its inclusion probability in the uniform exact layer.

\begin{lemma}[Core forest pattern bound]
Assume $j<g(K_v)$. For $B\subseteq A$ with $s$ clauses, if $B$ can occur inside the core then its support is a hyperforest. If that forest has $c$ components,
\[\Pr[B\subseteq\mathbf K]\le\frac{v^c(2D)^s}{(q)_{2s+c}}4^{-s},\qquad D=9.\]
\end{lemma}
\begin{proof}
Root each forest component at one vertex. Choose its image in at most $v$ ways. Order the hyperedges away from the root. Every new edge has exactly one previously embedded vertex, at most $D$ incident support-edge choices, and at most two bijections for its two new labeled vertices. This gives at most $v^c(2D)^s$ support embeddings. A fixed inverse embedding occurs under the random injection with probability $(q)_{2s+c}^{-1}$. Forest independence gives the exact factor $4^{-s}$ for the requested NAE signs.
\end{proof}

\begin{theorem}[Uniform relative moment matching]
Uniformly for every signed clause pattern $A$ of size $j=O(\log q)$ with $j<g(K_v)$,
\[\boxed{\left|\frac{\Pr[A\subseteq Y]}{u_j}-1\right|\le q^{\alpha-1+o(1)}\left(1+\frac D4\right)^j.}\]
In particular, every fixed-order inclusion moment matches the uniform exact layer to relative error $O_j(q^{\alpha-1})$.
\end{theorem}
\begin{proof}
Let $S=|A\cap\mathbf K|$. Conditional on the realized core, the filler gives
\[\Pr[A\subseteq Y\mid\mathbf K]=a_S,\qquad a_s=\frac{(t-k)_{j-s}}{(M-k)_{j-s}}.\]
Use the telescoping identity
\[\mathbb E a_S=a_0+\sum_{s=1}^j(a_s-a_{s-1})\Pr[S\ge s].\]
The no-core term satisfies $a_0/u_j=1+O(jk/t)=1+O(jv/q)$. For $s\ge1$, union-bound over the $s$-subsets $B\subseteq A$ and apply the forest pattern lemma. Since every nonempty forest has at least one component,
\[\Pr[S\ge s]\le q^{\alpha-1+o(1)}\binom js(D/2)^s q^{-2s}.\]
Also $a_s/u_j\le q^{2s}2^{-s+o(s)}$. Multiplying and summing over $s$ yields
\[q^{\alpha-1+o(1)}\sum_{s=1}^j\binom js(D/4)^s=q^{\alpha-1+o(1)}\big((1+D/4)^j-1\big),\]
which absorbs the no-core error.
\end{proof}

\subsection{A logarithmic moment window}
For the explicit $D=9$ support sequence, let
\[\gamma=\frac{2}{3\log11}.\]
The girth condition permits $j<c\log q$ whenever $c<\gamma\alpha$. The relative-moment error tends to zero whenever
\[c\log(13/4)<1-\alpha.\]
Optimizing the two constraints gives
\[\alpha_* = \frac{3\log11}{3\log11+2\log(13/4)}=0.75318699\ldots\]
and
\[\boxed{c_* = \frac{2}{3\log11+2\log(13/4)}=0.20940225\ldots.}\]

\begin{corollary}[Logarithmic-order pseudorandom UNSAT distribution]
For every constant $c<c_*$, along an infinite sequence of $q$ there is a distribution $\mu_q^+$ supported entirely on exact-density-$4/3$ UNSAT signed NAE formulas such that, uniformly for every set $A$ of at most $c\log q$ clause coordinates,
\[\boxed{\Pr_{\mu_q^+}[A\subseteq F]=(1+o(1))\frac{(t)_{|A|}}{(M)_{|A|}}.}\]
\end{corollary}

This closes the previous \emph{positive-distribution moment construction} obligation through a logarithmic order with an explicit constant. The remaining difficulty is now circuit-theoretic: prove that a polynomial-size monotone/set circuit can be approximated or tested using only moments inside this window, with accumulated error $o(1)$.

\subsection{Hardness and fixed-girth compatibility}
Independently of the planted construction, Kun's sparse incomparability theorem gives a polynomial reduction from NAE-CSP to NAE-CSP of any prescribed fixed girth, preserving satisfiability. Thus for every fixed $r$, NAE remains hard on supports of girth $>r$. Combining this with the gauge theorem shows that computational hardness and exact $r$-wise sign independence are compatible for every fixed $r$.

This does not prove hardness at girth $\Theta(\log q)$: Kun's formal theorem fixes the girth parameter. The author's later note that his expander construction yields $\log\log n$ girth is not promoted here to a growing-girth NP-hardness theorem without a separate uniform parameter audit. The logarithmic-girth positive distribution above needs no such hardness claim because its members are unconditionally UNSAT.

\section{TESTS AND OUTCOME}
The finite audit exhaustively enumerated every nonempty 3-uniform hypergraph on $n=3,4,5$ vertices, 1,039 hypergraphs total. For each support it built the two-bit-per-edge gauge matrix and verified
\[\operatorname{rank}\delta_J=v_J-c_J,\qquad 2|E|-\operatorname{rank}\delta_J=2|E|-v_J+c_J.\]
There were 2,078 assertions and zero discrepancies. The observed cycle-defect histogram ranged from zero (30 forest cases) through 16. Script SHA-256 is
\[\texttt{25476eb999d101a016c29f1dbb66834c2faa6c9dcdc61ca389650b554f0b325b}.\]
Raw-result SHA-256 is
\[\texttt{be1e99631180fbc76e3ed7874a9efa5b913cdf42b700698fe2bb759f295e147e}.\]
The computation is \COMP and only audits the symbolic theorem.

\paragraph{What changed.}
FTNCNAESLICE1 supplied a clean exact hard NAE slice and even-Fourier HD hierarchy. SLICEOUTPUTMOMENT1 left the asymptotic moment-matching UNSAT distribution open. The present report supplies such a distribution through $\Theta(\log q)$ moment order. The remaining load-bearing arrow is no longer positive-distribution construction; it is a circuit-to-low-moment or mass-production theorem.

\paragraph{Next target.}
Insert the logarithmic moment window into the Bonferroni overlap framework of SLICEBONFERRONI1. Determine whether robust-sunflower closure can reduce every size-$M^K$ monotone circuit to a term family whose required moment order is below $c_*\log q$ and whose factorial overlap tail is $o(1)$. Because the 2026 derandomized sunflower theorem only needs $O(w)$-wise independence for its own sunflower conclusion, the exact gauge-code theorem should be used to eliminate sign correlations while the remaining argument controls support overlap and DAG mass production.

Root \texttt{main.tex} remains truncated and was not modified.

\begin{thebibliography}{5}
\bibitem{Kun} G. Kun, \emph{Constraints, MMSNP and expander relational structures}, Combinatorica 33:475--486, 2013; arXiv:0706.1701.
\bibitem{EL} D. Ellis and N. Linial, \emph{On Regular Hypergraphs of High Girth}, Electron. J. Combin. 21(1):P1.54, 2014.
\bibitem{DD} A. Darmann and J. D\"ocker, \emph{On a Simple Hard Variant of Not-All-Equal 3-Sat}, Theor. Comput. Sci. 815:147--152, 2020.
\bibitem{Sun26} B. P. Cavalar, T. Fabris, P. Mukhopadhyay, S. Srinivasan, A. Yehudayoff, \emph{Derandomized Sunflowers and Radical Lower Bounds}, ECCC TR26-220, 2026.
\bibitem{Valiant} L. G. Valiant, \emph{Negation Is Powerless for Boolean Slice Functions}, SIAM J. Comput. 15(2):531--535, 1986.
\end{thebibliography}
\end{document}
% FINAL_VERIFICATION_SENTINEL::MAIN::NAEGAUGEHIGHGIRTH1::2026-10-17_1800::P_VS_NP_OPEN